\subsection{A Basis Is a Pair of Independent Directions}

Up to this point every vector has been described with respect to the horizontal
and vertical coordinate directions. That choice was natural because the
Cartesian axes had already been constructed, but it was still only a choice.
The geometric vector itself does not contain a preferred pair of coordinate
axes.

\begin{corebox}[The question of basis coordinates]
What numbers describe the same vector when we choose two different coordinate
directions? The vector remains fixed. Only the language used to describe it
changes.
\end{corebox}

\begin{definitionbox}
Two nonparallel vectors $\mathbf b_1$ and $\mathbf b_2$ form a basis of the
plane. Every vector $\mathbf u$ can be written in one and only one way as
\[
\boxed{\mathbf u=c_1\mathbf b_1+c_2\mathbf b_2}.
\]
The numbers $c_1$ and $c_2$ are the coordinates of $\mathbf u$ in the ordered
basis
\[
\mathcal B=(\mathbf b_1,\mathbf b_2),
\qquad
[\mathbf u]_{\mathcal B}=(c_1,c_2).
\]
\end{definitionbox}

The basis vectors need not be perpendicular and need not have length one. They
must only be nonparallel. Otherwise both directions collapse onto one line and
cannot describe every displacement in the plane.

\begin{interpretationbox}[Coordinates as signed instructions]
The equation
\[
\mathbf u=c_1\mathbf b_1+c_2\mathbf b_2
\]
means: move by $c_1\mathbf b_1$, then by $c_2\mathbf b_2$. The endpoint must
coincide with the endpoint of $\mathbf u$. The coefficients are signed amounts
of the chosen basis directions, not lengths of perpendicular shadows.
\end{interpretationbox}

\begin{center}
\begin{tikzpicture}[scale=0.95,>=stealth]
    \coordinate (O) at (0,0);
    \coordinate (B1) at (2.0,0.8);
    \coordinate (B2) at (0.8,1.8);
    \coordinate (TwoB1) at (4.0,1.6);
    \coordinate (U) at (4.8,3.4);

    \foreach \i in {-1,0,1,2,3} {
        \draw[gray!45,dashed]
            ({-0.8+2.0*\i},{-1.8+0.8*\i}) --
            ({2.4+2.0*\i},{5.4+0.8*\i});
    }
    \foreach \j in {-1,0,1,2,3} {
        \draw[gray!45,dashed]
            ({-2.0+0.8*\j},{-0.8+1.8*\j}) --
            ({6.0+0.8*\j},{2.4+1.8*\j});
    }

    \draw[blue!70!black,line width=1.1pt,->] (O) -- (B1)
        node[midway,below] {$\mathbf b_1$};
    \draw[orange!85!black,line width=1.1pt,->] (O) -- (B2)
        node[midway,left] {$\mathbf b_2$};
    \draw[blue!70!black,line width=1.1pt,->] (O) -- (TwoB1)
        node[midway,below right] {$2\mathbf b_1$};
    \draw[orange!85!black,line width=1.1pt,->] (TwoB1) -- (U)
        node[midway,right] {$\mathbf b_2$};
    \draw[very thick,->] (O) -- (U)
        node[midway,above left] {$\mathbf u$};
    \draw[blue!70!black,dashed] (B2) -- (U);
    \draw[orange!85!black,dashed] (TwoB1) -- (U);
    \fill (O) circle (2pt) node[below left] {$O$};
    \fill (U) circle (2pt);
    \node[align=left,anchor=west] at (5.3,3.0)
        {$\mathbf u=2\mathbf b_1+\mathbf b_2$\\
         $[\mathbf u]_{\mathcal B}=(2,1)$};
\end{tikzpicture}
\end{center}
